\documentclass[11pt,a4paper]{article}
\usepackage[utf8]{inputenc}
\usepackage{amsmath,amssymb,amsthm}
\usepackage{hyperref}
\usepackage[margin=1in]{geometry}

\title{HRAC: History-Residual Adaptive Clipping for Byzantine-Robust Federated Aggregation}
\author{For reviewing}
\date{}

\begin{document}
\maketitle


\section{Implementation}

\subsection{Global norm cap (MSA-style defense)}

At round $t$, each client $i$ sends a local update $\Delta_i^{(t)} \in \mathbb{R}^d$. Let $n_i = \|\Delta_i^{(t)}\|$ and $\bar{n} = \mathrm{median}(n_1,\ldots,n_N)$. Define the global cap
\begin{equation}
  B = \max\bigl(c_g \cdot \bar{n},\; \epsilon\bigr), \qquad
  s_i = \min\left(1,\; \frac{B}{n_i + \epsilon}\right).
\end{equation}
The server forms globally capped messages $\widetilde{\Delta}_i = s_i \cdot \Delta_i^{(t)}$. Only scaling is used; large norms are shrunk to at most $B$, small norms are unchanged. This prevents a single client from controlling the scale of the round. Typical choice: $c_g = 3$.

\subsection{Small-norm lift (for statistics only)}

If any client has $n_i < 0.5\,\bar{n}$, we define an ``effective'' message for \emph{statistics} (used to compute $\tau,\mu,\nu$ and weights), while \emph{aggregation} still uses the real capped update. For client $i$ with $n_i < 0.5\,\bar{n}$:
\begin{equation}
  \Delta_i^{\mathrm{eff}} = \frac{\bar{n}}{n_i + \epsilon}\,\Delta_i^{(t)} \quad \Rightarrow \quad \|\Delta_i^{\mathrm{eff}}\| = \bar{n}.
\end{equation}
For others, $\Delta_i^{\mathrm{eff}} = \widetilde{\Delta}_i$. So small-norm clients do not get scaled updates in the aggregate; only the quantities that drive $\tau,\mu,\nu$ are computed from the lifted scale so that statistics are comparable across clients.

\subsection{Per-client residual clipping and processed update}

The server keeps per-client state: $\mathbf{b}_i$ (aggregation baseline), $\mathbf{b}_i^{\mathrm{norm}}$ (stats baseline), $\mu_i$ (norm scale), $\nu_i$ (change scale), $\mathbf{r}_i^{\mathrm{prev}}$ (previous clipped residual for stats). The clipping threshold is $\tau_i = c \cdot \mu_i + \epsilon$ with $c = 2.5$.

For aggregation, the server uses the real baseline $\mathbf{b}_i$ and capped message $\widetilde{\Delta}_i$:
\begin{equation}
  \mathbf{r}_i = \widetilde{\Delta}_i - \mathbf{b}_i, \qquad
  \bar{\mathbf{r}}_i = \mathrm{clip}(\mathbf{r}_i,\; \tau_i), \qquad
  \widehat{\Delta}_i = \mathbf{b}_i + \bar{\mathbf{r}}_i.
\end{equation}
Here $\mathrm{clip}(\mathbf{r},\tau)$ scales $\mathbf{r}$ so that its norm does not exceed $\tau$ (scale factor $\min(1, \tau/\|\mathbf{r}\|)$). Then $\widehat{\Delta}_i$ is re-scaled so that $\|\widehat{\Delta}_i\| \le B$ (again by scaling only). The result is the processed update $\Delta_i^{\mathrm{proc}}$ used in aggregation.

For the stats path (used to update $\mu_i,\nu_i,\mathbf{b}_i^{\mathrm{norm}},\mathbf{r}_i^{\mathrm{prev}}$): when the small-norm lift is active, the same formulas are applied with $\Delta_i^{\mathrm{eff}}$ and $\mathbf{b}_i^{\mathrm{norm}}$ instead of $\widetilde{\Delta}_i$ and $\mathbf{b}_i$, yielding $\bar{\mathbf{r}}_i^{\mathrm{norm}}$ and a normalized $\Delta_i^{\mathrm{norm}}$; otherwise the stats path uses the same $\bar{\mathbf{r}}_i$ and $\widehat{\Delta}_i$ as above.

\subsection{Weights (nu-penalty)}
\label{sec:weights}

Weights are applied only when $t > t_0$ (e.g.\ $t_0 = 50$). Let $\bar{\nu} = \frac{1}{N}\sum_{i=1}^N \nu_i$ (mean of $\nu$ over clients). We penalize only clients with $\nu_i > 1$ using the excess above the mean:
\begin{equation}
  e_i = \max\bigl(0,\; \nu_i - \bar{\nu}\bigr) \quad \text{if }\nu_i > 1;\qquad e_i = 0 \quad \text{otherwise}.
\end{equation}
With $\alpha = 5$,
\begin{equation}
  \tilde{w}_i = \frac{1}{1 + \alpha \cdot e_i}, \qquad
  w_i = \frac{\tilde{w}_i}{\sum_j \tilde{w}_j}.
\end{equation}
So clients whose $\nu$ is far above the mean get smaller weights; using $\bar{\nu}$ (mean) instead of the median keeps penalizing outliers even when the mean is inflated (e.g.\ many clients with $\nu > 1$). Then an iterative cap is applied so that no single weight exceeds $w_{\max} = 0.3$, and weights are renormalized after each cap.

\subsection{Aggregation and optional momentum (bypass only)}

The aggregate update is $\mathbf{g}^{(t)} = \sum_{i=1}^{N} w_i\,\Delta_i^{\mathrm{proc}}$, with weights $w_i$ from the $\nu$-penalty (Section~\ref{sec:weights}). In the \emph{normal} path, the server updates the model using $\mathbf{g}^{(t)}$ as the step direction (the aggregator returns $\mathbf{g}^{(t)}$).

An \emph{aggregate momentum} $\mathbf{m}^{(t)}$ is maintained internally:
\begin{equation}
  \mathbf{m}^{(0)} = \mathbf{g}^{(0)}, \qquad
  \mathbf{m}^{(t)} = \beta\,\mathbf{m}^{(t-1)} + (1-\beta)\,\mathbf{g}^{(t)} \quad (t \ge 1),
\end{equation}
with $\beta \in (0,1)$ (e.g.\ $\beta = 0.95$). This momentum is \emph{not} used for the model update in the normal path. It is used only in the \emph{bypass} path (e.g.\ when a scale anomaly is detected): in that case the server updates the model using $\mathbf{m}^{(t)}$ instead of $\mathbf{g}^{(t)}$, so that a single bad round is dampened by past rounds. Optionally, a norm cap on $\mathbf{g}^{(t)}$ can be applied (e.g.\ cap $= \min(\lambda B,\, k \cdot \mathrm{median}(\|\mathbf{g}^{(1)}\|,\ldots)$); in the implementation this can be disabled.


\subsection{State update}

Let $\Delta_i^{\mathrm{proc},\mathrm{cap}}$ denote the processed update for client $i$ after the final scaling to stay within $B$. For the stats path, $\|\bar{\mathbf{r}}_i^{\mathrm{norm}}\|_{\mathrm{eff}}$ is $\|\Delta_i^{\mathrm{norm}}\|$ for small-norm clients and $\|\bar{\mathbf{r}}_i^{\mathrm{norm}}\|$ otherwise. The change term is $d_i = \|\bar{\mathbf{r}}_i^{\mathrm{norm}} - \mathbf{r}_i^{\mathrm{prev}}\|$, optionally capped (e.g.\ $d_i^{\mathrm{capped}} = \min(d_i, \nu_i \cdot \gamma)$ with $\gamma$ a constant or $\mathrm{None}$ to disable). The server updates:
\begin{align}
  \mathbf{b}_i &\leftarrow \rho_b\,\mathbf{b}_i + (1-\rho_b)\,\Delta_i^{\mathrm{proc},\mathrm{cap}}, \\
  \mathbf{b}_i^{\mathrm{norm}} &\leftarrow \rho_b\,\mathbf{b}_i^{\mathrm{norm}} + (1-\rho_b)\,\Delta_i^{\mathrm{norm}}, \\
  \mu_i &\leftarrow \max\bigl(\rho_\mu\,\mu_i + (1-\rho_\mu)\,\|\bar{\mathbf{r}}_i^{\mathrm{norm}}\|_{\mathrm{eff}},\; \mu_{\min}\bigr), \\
  \nu_i &\leftarrow \max\bigl(\rho_\nu\,\nu_i + (1-\rho_\nu)\,d_i^{\mathrm{capped}},\; \nu_{\min}\bigr), \\
  \mathbf{r}_i^{\mathrm{prev}} &\leftarrow \bar{\mathbf{r}}_i^{\mathrm{norm}}.
\end{align}
A single smoothing coefficient $\rho_\nu$ is used for $\nu_i$ (no switching by aggregate momentum or $d_i < \nu_i$). Typical parameters: $\rho_b = 0.98$, $\rho_\mu = 0.95$, $\rho_\nu = 0.87$ or $0.95$; $\mu_{\min},\nu_{\min}$ are small constants to prevent collapse.

\end{document}
